\documentclass[11pt]{article}
\usepackage[margin=1in]{geometry}
\usepackage[T1]{fontenc}
\usepackage[utf8]{inputenc}
\usepackage{amsmath,amssymb,amsthm}
\usepackage{enumitem}

\title{ICPC World Finals 2024\\A. Billboards}
\author{}
\date{}

\begin{document}
\maketitle

\section*{Problem Summary}

Each sponsor has a nonnegative continuous value density along the billboard. We must partition the
billboard into contiguous pieces, assign one piece to each sponsor, and guarantee that sponsor $i$
receives value at least $1/n$ of its total value.

\section*{Key Observations}

\begin{itemize}[leftmargin=*]
    \item Let $\mu_i([a,b])$ be sponsor $i$'s value on interval $[a,b]$. Because the density is continuous
    and nonnegative, $\mu_i([x,y])$ is continuous in $y$ and never decreases.
    \item Fix a current left boundary $x$. If sponsor $i$ cuts the \emph{leftmost} point $c_i \ge x$ such
    that $\mu_i([x,c_i]) = \mu_i([0,\ell])/n$, then every point before $c_i$ is still worth at most that
    amount to sponsor $i$.
    \item Therefore, if we choose the sponsor whose leftmost such cut is smallest, every other remaining
    sponsor loses at most $1/n$ of its total value from the left prefix we remove.
    \item This is exactly the invariant needed for a proportional contiguous allocation: after removing one
    piece, each remaining sponsor still values the suffix enough for the remaining number of pieces.
\end{itemize}

\section*{Algorithm}

\begin{enumerate}[leftmargin=*]
    \item For each sponsor, precompute the prefix integral of its piecewise linear density at every breakpoint.
    This gives the total value $T_i$ and the target value $T_i / n$.
    \item Maintain the current left boundary $L$ and the set of sponsors not assigned yet.
    \item While more than one sponsor remains:
    \begin{enumerate}[leftmargin=*]
        \item For every remaining sponsor $i$, find the smallest point $c_i \ge L$ such that
        $\mu_i([L,c_i]) = T_i / n$.
        \item Pick the sponsor $p$ with minimum $c_i$.
        \item Assign interval $[L,c_p]$ to sponsor $p$, output $(c_p, p)$, and set $L = c_p$.
    \end{enumerate}
    \item Assign the remaining suffix $[L,\ell]$ to the last sponsor.
\end{enumerate}

To evaluate $\mu_i([0,x])$ quickly, we binary search the segment containing $x$ and integrate the linear
function on that segment. To recover the leftmost cut point from a target prefix area, we binary search the
first breakpoint whose prefix area reaches that target, then solve one linear or quadratic equation inside that
segment.

\section*{Correctness Proof}

We prove that the algorithm always outputs a valid allocation.

\paragraph{Lemma 1.}
Suppose the current left boundary is $L$, sponsor $i$ is still unassigned, and
$\mu_i([L,\ell]) \ge k \cdot T_i / n$ where $k$ is the number of unassigned sponsors. Then there exists a
smallest point $c_i \in [L,\ell]$ such that $\mu_i([L,c_i]) = T_i / n$.

\paragraph{Proof.}
At $y=L$ the value $\mu_i([L,y])$ is $0$, and at $y=\ell$ it is at least $T_i/n$. The function
$\mu_i([L,y])$ is continuous and nondecreasing in $y$, so by the intermediate value theorem there exists
at least one such point. Taking the smallest one is well-defined. \qed

\paragraph{Lemma 2.}
When the algorithm assigns the next piece, every still-unassigned sponsor loses value at most $T_i / n$
from the removed prefix.

\paragraph{Proof.}
Let $p$ be the sponsor selected by the algorithm, so $c_p \le c_j$ for every other remaining sponsor $j$.
By definition, $c_j$ is the \emph{smallest} point where sponsor $j$ has accumulated value exactly $T_j/n$
from the current left boundary. Since $c_p \le c_j$, the prefix $[L,c_p]$ is entirely before that first
threshold for sponsor $j$, hence $\mu_j([L,c_p]) \le T_j/n$. \qed

\paragraph{Lemma 3.}
After each assignment, if $k-1$ sponsors remain, then every remaining sponsor $j$ values the remaining
suffix at least $(k-1) \cdot T_j / n$.

\paragraph{Proof.}
Before the step, the inductive hypothesis gives $\mu_j([L,\ell]) \ge k \cdot T_j / n$. By Lemma 2,
the removed interval is worth at most $T_j/n$ to sponsor $j$. Therefore
\[
\mu_j([c_p,\ell]) = \mu_j([L,\ell]) - \mu_j([L,c_p])
\ge k \cdot T_j / n - T_j / n
= (k-1) \cdot T_j / n.
\]
\qed

\paragraph{Theorem.}
The algorithm outputs a contiguous allocation where every sponsor receives at least $T_i / n$ value.

\paragraph{Proof.}
Initially $k=n$ and every sponsor trivially values the whole billboard at $n \cdot T_i / n = T_i$, so the
invariant of Lemma 3 holds. By Lemma 1 the next cut always exists, and by Lemma 3 the invariant is
preserved after each assignment.

When only one sponsor remains, the invariant says that this sponsor values the remaining suffix at least
$1 \cdot T_i / n$, so assigning the entire suffix is valid. All assigned pieces are contiguous, disjoint, and
their union is the whole billboard because each step cuts off a prefix of the remaining suffix. Thus every
sponsor receives a valid piece of value at least $T_i / n$. \qed

\section*{Complexity Analysis}

Let $M$ be the total number of breakpoints over all sponsors. Precomputing prefix areas costs $O(M)$.

For each of the first $n-1$ assignments, we scan all remaining sponsors. For one sponsor we do two binary
searches over its breakpoints: one to evaluate the prefix integral at the current left boundary, and one to
locate the segment containing the target cut. Thus the total running time is
\[
O\!\left(M + \sum_{i=1}^{n} n \log m_i \right),
\]
which is safely within the limits for $n \le 5000$ and $\sum m_i \le 500000$. The memory usage is $O(M)$.

\section*{Implementation Notes}

\begin{itemize}[leftmargin=*]
    \item Prefix areas are stored as \texttt{long double} because cut points are real numbers.
    \item Inside a segment, recovering the cut point means solving
    $b \Delta x + \frac{s}{2} \Delta x^2 = \text{needed area}$, where $b$ is the left endpoint value and
    $s$ is the slope on that segment.
    \item The statement allows tiny absolute or relative error, so printing with 15 digits after the decimal
    point is more than enough.
\end{itemize}

\end{document}
